\selectlanguage{russian}
% TODO: the Russian abstract is a machine-drafted translation of the English one -- please have it checked.
Большие визуально-языковые модели (VLM) часто дают гладкие ответы, не подтверждаемые изображением. В работе рассматривается детальная постановка задачи из соревнования SHROOM-Visions 2026: по изображению, вопросу и ответу модели на английском, французском, итальянском или китайском языке требуется отметить каждый галлюцинированный фрагмент текста с точностью до символа, отнести его к одной из пяти категорий (выдумка, неверное описание, ошибка распознавания текста, ошибка подсчёта, прочее) и указать вероятность, отражающую долю аннотаторов, которые отметили бы этот фрагмент.

Предложенный детектор состоит из четырёх частей. (i) Модель-судья Qwen3-VL-8B, дообученная методом QLoRA, выписывает галлюцинированные фрагменты в формате JSON. (ii) Декодер фрагментов, учитывающий метки, переводит этот текст в позиции символов и сохраняет фрагменты, которым аннотаторы дали две метки. (iii) Модуль самосогласованности объединяет восемь сэмплированных ответов и жадный ответ судьи посимвольным голосованием отдельно по каждой метке. (iv) Зонд неопределённости типа «белого ящика» читает энтропию токенов, правдоподобие, долю внимания к контексту и скрытое состояние замороженной модели и декодирует фрагменты с помощью кодировщика BiGRU и пяти совместно обучаемых бинарных CRF; он участвует в голосовании как ещё один голосующий. Итоговый ответ сохраняет вероятность для каждого символа, а не одно значение на фрагмент.

Все решения принимались по заранее зарегистрированному протоколу. Валидационная выборка разделена на часть для выбора и часть для подтверждения, сгруппированные по изображениям. В каждом решении сравнивалось не более пяти кандидатов. Каждое утверждение требовало трёх случайных инициализаций и парного бутстрепа. В среднем по трём инициализациям на официальной тестовой выборке детектор достигает Cor+Lbl 0,4296 против 0,4132 у той же модели, дообученной только на исходных данных (+0,0163). Улучшение наблюдается для каждой инициализации и каждого языка; IoU растёт на 0,018, Cor — на 0,014. Такой результат соответствует 5-му, 5-му, 5-му и 3-му местам в рейтингах для английского, французского, итальянского и китайского языков (наша заявка на соревнование занимала 8-е, 6-е, 7-е и 5-е места).

Отдельная серия экспериментов изолированно проверяла синтетические обучающие данные. Ни одно из пяти условий не превзошло базовую модель, дообученную на исходных данных, в том числе перенос реальных фрагментов, которые текстовый классификатор не отличает от настоящих. Разложение метрики на вклад «чистых» и размеченных ответов объясняет причину: синтетические примеры повышают долю галлюцинированных примеров в обучении, и выигрыш на ответах с галлюцинациями теряется из-за ложных срабатываний на чистых ответах.
